\subsection{Related work}\label{app:other_works}

Here, we discuss the differences to other works on AI safety and one related position paper that was concurrently submitted to our work.

\subsubsection{Difference to other works related to AI safety}

The development of safe and reliable AI, particularly Large Language Models (LLMs), is a central concern. As a result, there are numerous works which address the importance of reliable and comprehensive safety evaluations in the LLM safety community~\cite{bommasani2021opportunities, weidinger2021ethical, ganguli2022red, liang2022holistic, hendrycks2023overview, srivastava2022beyond, chen2024trustworthy, bai2022constitutional, shevlane2023model, bowman2022measuring, sillberg2024eu, casper2023open, anwar2024foundational, casper2024black, schwinn2023adversarial, casper2023explore, reuel2024open}.

Yet, all of these contributions have substantially different goals compared to the proposed position. One group of works primarily evaluate capabilities and outline potential risks~\cite{bommasani2021opportunities, weidinger2021ethical, hendrycks2023overview, srivastava2022beyond, liang2022holistic, shevlane2023model, anwar2024foundational}, another is focused on methodology and technical limitations and advancements~\cite{ganguli2022red, chen2024trustworthy, bai2022constitutional, bowman2022measuring, casper2023open, casper2023explore}, while others focus on formulating policy recommendations required for safe AI usage now and in the future~\cite{sillberg2024eu, casper2024black, reuel2024open}. Our work, in contrast, focuses specifically on the process of research itself within adversarial alignment, diagnosing how current practices risk hindering cumulative progress by not sufficiently learning from the history of adversarial robustness research.

In the following we provide a categorization of prior work from the AI safety community. Note, that some work could be assigned to multiple categories.

\textbf{Discussions on AI capabilities and risks.} A significant body of work explores the capabilities of LLMs and the spectrum of risks they introduce. For instance, \citet{bommasani2021opportunities} provide a broad overview of foundation models capabilities and resulting risks, while \citet{weidinger2021ethical} detail potential ethical and social harms stemming from language models. The potential for catastrophic risks is further analyzed in \citet{hendrycks2023overview, shevlane2023model}. Other research highlights areas of concern giving potential future capabilities of LLMs and discusses current limitations \cite{srivastava2022beyond}. More recently, foundational challenges in ensuring LLM alignment and safety have been systematically documented, pointing to the deep-seated difficulties in achieving these goals \cite{anwar2024foundational}.

\textbf{Methodological improvements for AI safety.} The LLM safety community has proposed numerous methodological advancements. This includes novel alignment techniques, such as Constitutional AI, which aims to instill harmlessness through AI feedback \cite{bai2022constitutional}. Red teaming has emerged as a critical methodology for discovering vulnerabilities \cite{ganguli2022red, casper2023explore}, and significant effort is dedicated to measuring progress on scalable oversight \cite{bowman2022measuring}. Methodological critiques also exist, for example, highlighting the limitations of RLHF \cite{casper2023open}. 

\textbf{AI policy and governance recommendations.} Recognizing the societal impact of AI, researchers and institutions are increasingly proposing policy and governance frameworks. These range from critiques and recommendations concerning specific regulations like the EU AI Act \cite{sillberg2024eu}, to broader architectural frameworks for trustworthy and responsible AI \cite{chen2024trustworthy}, and discussions on open problems in technical AI governance \cite{reuel2024open}. Such works aim to guide the development and deployment of AI in a manner that aligns with societal values and minimizes harm. Moreover, they advocate for safety policy, where technical solutions alone might fail. 

In summary, existing literature offers critical insights into AI risks, develops novel methodologies, and proposes governance structures. Our work complements these efforts by adopting a historical and methodological lens focused on how to \textbf{facilitate research progress}. By drawing explicit lessons from the trajectory of adversarial robustness research in other domains, we identify current pitfalls and advocate for a realignment of research objectives to ensure that the field makes efficient and durable progress, avoiding the cycles of ineffective solutions that characterized earlier efforts in adversarial machine learning~\cite{madry_towards_2018, athalye_obfuscated_2018, tramer_adaptive_2020, schwinn2023adversarial}. This specific focus distinguishes our contribution.

\subsubsection{Concurrent position}

Our findings are reinforced by the concurrent work of~\citet{rando2025adversarial}, who independently identify similar fundamental challenges in adversarial alignment of LLMs. 
Their research offers detailed case studies of specific threat models like jailbreaks, poisoning, and unlearning. We also investigate specific failures as--albeit in less detail. Moreover, in contrast to their work, we develope a high-level taxonomy encompassing all sorts of threat models. Furthermore, we go beyond problem identification to propose actionable, forward-looking solutions that address emerging challenges in adversarial alignment research.


% %Their paper offers detailed case studies of specific threat models like jailbreaks, poisoning, and unlearning, whereas our work develops a high-level taxonomy and focuses more on forward-looking solutions. 
% Together, both works offer complementary perspectives toward a more comprehensive understanding of the identified issues.


\subsection{Further arguments against the position}~\label{app:arguments_against}

Another argument against a focus on well-defined sub-problems is the broad attack surface of foundation models from chat assistants to autonomous agents. The majority of threats related to unexplored and upcoming applications will not be formally well-defined. Here, research requires identifying novel threat surfaces in new application areas. This kind of research has been proven to be highly relevant in the past, including the identification of vulnerabilities in commercial hashing algorithms~\cite{levenson2021apple}, seemingly reliable protections of art against generative AI~\cite{honig2024adversarial}, or vulnerabilities of newly developed benchmarks, such as LLM leaderboard~\cite{huang2025exploring}.

\textbf{New opportunities for defenses.} In this work, we primarily emphasize the increased complexity of achieving robustness in LLMs compared to previous research domains. However, the semantic output and input space of LLMs not only complicate measuring robustness but also open up new opportunities for defenses. Early adversarial detection methods in computer vision focused on identifying \textit{small}, \textit{imperceptible} norm-bounded perturbations, proved largely ineffective~\cite{carlini2017adversarial}. Moreover, detection in the output space was not possible due to the prevalence of simple classification tasks. Conversely, recent advancements in the LLM domain demonstrate the strong empirical results of input and output classifiers for defending against harmful requests~\cite{kim2023robust, sharma2025constitutional}. This suggests that the inherent semantics of the LLMs domain, while presenting new challenges, also facilitate the development of new defense mechanisms previously unavailable in other domains.

\textbf{Capability progress is too fast.} Another argument against the proposed position of reliable research in constrained settings is the rapid pace of progress in model capabilities. By simplifying the problem, researchers risk creating solutions that will quickly become obsolete as AI models evolve. This could be specifically the case as current adversarial attack algorithms often rely on LLMs to jailbreak other models. Rather than trying to solve isolated, small problems in the field, researchers should concentrate on large, fundamental questions that are not dependent on current technical limitations, as they may be resolved with progress in LLM capability. Moreover, while open-source models are increasing in their capabilities, research on open-source models may always lag behind proprietary ones and academia may recover mistakes already identified by industry labs. 

\subsection{Beyond Discussed Threats}~\label{app:beyond}
Most arguments presented in this work regarding emerging challenges in LLM robustness research are applicable beyond the specific threat models discussed as part of the taxonomy. 
These include issues related to complex input-output domains (e.g., precisely defining attack constraints or evaluation goals) and challenges associated with large model sizes, vast number of hyperparameters, and varied implementation settings.
This includes integrity threats, such as model poisoning, as well as attacks on confidentiality (e.g., membership inference, model inversion) and hybrid threats, like model unlearning. 
However, we note that these related areas are likely to present their own unique complexities and challenges within the LLM context that are currently hindering research progress, which requires discussion in future work.

\subsection{Examples of Measurability Improvements in Recent Work}\label{app:measurability}
Here, we highlight recent datasets that demonstrate how complex problems can be made more measurable and list desirable properties they fulfill:
\begin{itemize}
    \item \textbf{Jailbreak Tax}~\cite{nikolic2025jailbreak}: Multiple choice outcomes, verifiable outcomes.
    \item \textbf{AutoAdvExBench}~\cite{carlini2025autoadvexbench}: Quantifiable metrics, real-world coding task.
    \item \textbf{NYU CTF Bench}~\cite{shao2025nyu}: Unit-testable, cyber security challenges.
    \item \textbf{AgentHarm}~\cite{andriushchenko2024agentharm}: Verifiable actions.
\end{itemize}

Beyond highlighting desirable properties in existing benchmarks, we propose three ideas to improve measurability:
\begin{enumerate}
    \item \textbf{Robust judgeability}: For open-ended datasets such as JailbreakBench or HarmBench, curate subsets of ``robustly judgeable examples'' with low judge disagreement across defenses, models, and attacks.
    \item \textbf{Direct dataset construction}: Directly construct datasets with straightforward success criteria, such as highly toxic completions identifiable by keywords.
    \item \textbf{Guaranteed success settings}: Define settings where an attack is guaranteed to succeed in principle. Following the approach proposed in~\cite{carlini2023aligned}, we identify extremely unlikely output sequences for a model via brute-force search and optimize the attack to elicit them. This ensures that the adversarial optimization problem can be solved and studied in isolation.
\end{enumerate}

\subsection{Evidence for Claims}\label{app:evidence}
In this section, we provide a more detailed breakdown of the central claims presented in this position paper. For each claim, we synthesize evidence from the literature to support our position and propose concrete experimental directions to further validate our claims in the future.

\begin{enumerate}
    \item \textbf{Claim:} \textit{Complex, open-ended evaluations often result in measurement errors, whereas measurable sub-problems provide more reliable signals of progress.}
    
    \textbf{Evidence:} 
    Evaluating the success of an adversarial attack on a broad, open-ended task (e.g., "generate a harmful response") relies heavily on subjective interpretation and complex automated judges, which introduces significant noise~\cite{beyer2025llm}. Recent benchmarks demonstrate that narrowing the domain improves reliability. For instance, the \textbf{Jailbreak Tax}~\cite{nikolic2025jailbreak} reveals that open-ended evaluations often systematically overestimate attack success due to high false-positive rates in LLM judges. By restricting the evaluation to a "measurable sub-problem", such as a multiple-choice format or verifiable domains (e.g., math), researchers can isolate the performance of the attack algorithm from the vagaries of the evaluation pipeline. Similarly, \textbf{NYU CTF Bench}~\cite{shao2025nyu} and \textbf{AgentHarm}~\cite{andriushchenko2024agentharm} show that unit-testable challenges (e.g., capturing a flag or performing a specific function call) offer unambiguous success criteria, eliminating the need for subjective judgment and improving measurability.
    
    \textbf{Proposed Experiment:} 
    Future studies could correlate performance improvements on these measurable, narrow tasks with performance on broader, open-ended datasets. If improvements in optimization on narrow tasks consistently transfer to open-ended settings, it would validate the use of measurable sub-problems as reliable proxies for general progress.

    \item \textbf{Claim:} \textit{Open-ended evaluations suffer from significant variability in judge performance, undermining comparability.}
    
    \textbf{Evidence:} 
    The reliance on "LLM-as-a-Judge" for evaluating open-ended adversarial attacks introduces a critical source of instability. Studies such as those by \citet{li2025llms} and \citet{beyer2025llm} highlight substantial inconsistencies between different LLM judges, as well as discrepancies between LLM judges and human annotators. These judges can be sensitive to irrelevant formatting details or biases in the prompt, leading to false positives that vary across different attack methods. This variability means that a \enquote{successful} defense in one paper might fail under a slightly different judge configuration in another, making it impossible to reliably compare defense mechanisms.
    
    \textbf{Proposed Experiment:} 
    Extend existing studies to systematically investigate differences in false positive/negative rates of various judges specifically across different defense and attack algorithms. Quantifying how much the choice of judge alters the \textit{ranking} of defenses would demonstrate the severity of this issue.

    \item \textbf{Claim:} \textit{Reliance on proprietary model evaluations creates significant reproducibility risks.}
    
    \textbf{Evidence:} 
    Evaluations conducted on proprietary APIs (e.g., GPT, Claude) are inherently unstable due to unannounced updates, version deprecations, and opaque system changes. A defense or attack evaluated on "GPT-5.2" today may yield completely different results next week if the backend model or safety filter is updated. This \enquote{drifting baseline} makes it impossible to reproduce results from prior literature or to benchmark progress over time accurately. The deprecation of specific API snapshots further exacerbates this, rendering older experimental setups permanently inaccessible.
    
    \textbf{Proposed Experiment:} 
    Conduct a longitudinal study tracking the success rate of a fixed set of attacks against proprietary model APIs over a period of 6-12 months. Quantifying the variance in results solely due to API changes would provide concrete evidence of the \enquote{reproducibility crisis} in proprietary evaluations.

    \item \textbf{Claim:} \textit{Research progress relies on the adoption of common, standardized benchmarks.}
    
    \textbf{Evidence:} 
    The history of adversarial robustness in computer vision demonstrates the necessity of standardization. The creation of \textbf{RobustBench}~\cite{croce2020robustbench} provided a unified framework that allowed the community to rigorously track progress, filter out ineffective defenses, and identify genuine state-of-the-art methods. In contrast, the current state of LLM robustness research is fragmented, with many papers using custom datasets, evaluation protocols, and threat models~\cite{beyer2025llm}. This fragmentation obscures whether new methods are genuinely better or simply optimized for a specific, non-standard setting.
    
    \textbf{Proposed Experiment:} 
    Encourage and track the adoption of third-party evaluations for LLM defenses, similar to the community effort seen in RobustBench. A study analyzing the "breakage rate" of defenses when evaluated on a standardized benchmark versus their reported performance would highlight the gap caused by lack of standardization.

    \item \textbf{Claim:} \textit{The performance gap between open-source and proprietary models has narrowed, making open-source models viable proxies for research.}
    
    \textbf{Evidence:} 
    Historically, proprietary models were significantly more capable than open-source alternatives, justifying their use in evaluations. However, recent trends show a rapid convergence. Leaderboards like the \textbf{LMSYS Chatbot Arena} show open-weights models (e.g., GLM-4.7, Kimi-k2-Thinking-Turbo, and DeepSeek-V3.2) performing on par with or even exceeding frontier proprietary systems. This suggests that insights gained from robustifying or attacking these strong open models are likely to transfer to proprietary counterparts, without the reproducibility downsides of API-based research. 
    
    \textbf{Proposed Experiment:} 
    Track whether methodological trends (e.g., susceptibility to specific attack vectors, effectiveness of defenses) observed on state-of-the-art open models are predictive of trends in proprietary models. If the correlation is high, it validates the shift towards open-source-first research.

    \item \textbf{Claim:} \textit{Academic research on open models is a driver of safety improvements in proprietary systems.}
    
    \textbf{Evidence:} 
    Despite the focus on proprietary models in some benchmarks, foundational methodological advances in AI safety often originate in academia using open models. Key examples include the \textbf{GCG} attack~\cite{zou2023universal}, the \textbf{HarmBench} framework~\cite{mazeika2024harmbench}, and early adversarial training algorithms~\cite{xhonneux2024efficient}. These methods, developed on open weights, have heavily influenced the safety reports and defense strategies of major proprietary labs (e.g., referenced in Gemini and LLaMA technical reports), demonstrating that open-model research directly impacts the wider ecosystem.
    
    \textbf{Proposed Experiment:} 
    Perform a citation and genealogy analysis of safety techniques deployed in major industrial models to quantify the proportion of core safety methodologies that originated from academic research on open systems.

    \item \textbf{Claim:} \textit{Current automated attacks overestimate robustness compared to human adaptive attacks.}
    
    \textbf{Evidence:} 
    Automated attack algorithms, while improving, still struggle with the complex optimization landscape of language. Evidence from recent competitions and studies~\cite{li2024llm, andriushchenko2024jailbreaking, nasr2025attacker} shows that manual \enquote{jailbreaking} or human-in-the-loop strategies consistently achieve higher Attack Success Rates (ASR) on robust models than fully automated baselines. This gap implies that relying solely on current automated attacks for evaluation can lead to a false sense of security, as the model may be vulnerable to a more creative or adaptive human adversary.
    
    \textbf{Proposed Experiment:} 
    Benchmark state-of-the-art automated attacks against a coordinated human red-teaming effort on the same set of defended models at the end of 2026. Quantifying the \enquote{automation gap} would emphasize the need for stronger adaptive attack algorithms (or proxies) before automated evaluation can be fully trusted.
\end{enumerate}
